% 数论（综述）
% license CCBY4
% type Wiki

本文根据 CC-BY-SA 协议转载翻译自维基百科\href{https://en.wikipedia.org/wiki/Number_theory}{相关文章}

\begin{figure}[ht]
\centering
\includegraphics[width=6cm]{./figures/ae528936f6668ed9.png}
\caption{} \label{fig_SHuLun_1}
\end{figure}
数论是纯数学的一个分支，主要致力于研究整数以及算术函数。数论学家研究素数，以及由整数构造的数学对象（例如有理数）的性质，或作为整数的推广而定义的对象（例如代数整数）。

整数既可以单独研究，也可以作为某些方程的解来研究（丢番图几何）。数论中的问题常常能够通过研究某些解析对象来理解，例如黎曼$\zeta$函数，它以某种方式编码了整数、素数或其他数论对象的性质（解析数论）。另外，还可以研究实数与有理数之间的关系，例如研究无理数如何用分数来逼近（丢番图逼近）。

数论与几何一样，是数学最古老的分支之一。数论的一个特点是，它处理的往往是非常容易理解却极难解决的命题。例如，费马大定理在最初提出后的 358 年才被证明，而哥德巴赫猜想自 18 世纪提出以来至今仍未解决。德国数学家卡尔·弗里德里希·高斯曾经说过：“数学是科学的女王，而数论是数学的女王。”\(^\text{[1]}\)数论一度被视为纯粹数学的典范，认为它在数学之外没有任何应用，直到 20 世纪 70 年代，人们才发现素数可被用作构建公钥密码算法的基础。
\subsection{历史}

数论是数学的一个分支，研究整数及其性质与关系。\(^\text{[2]}\)整数是一个集合，它在自然数集合${1,2,3,\dots}$的基础上扩展，加入了数0，以及自然数的相反数${-1,-2,-3,\dots}$。数论学家不仅研究素数，还研究由整数构造的数学对象（例如有理数）的性质，或定义为整数推广的对象（例如代数整数）。\(^\text{[3][4]}\)

数论与算术密切相关，一些作者甚至将这两个词视为同义词。\(^\text{[5]}\)然而，今天“算术”一词通常指研究数值运算的学科，并扩展到实数范围。\(^\text{[6]}\)在更具体的意义上，数论仅限于研究整数，关注它们的性质和相互关系。\(^\text{[7]}\)传统上，它也被称为高等算术。\(^\text{[8]}\)到 20 世纪初，“数论”这一术语已被广泛采用。\(^\text{[note 1]}\)其中“number”意为整数，既可以指自然数，也可以指全体整数。\(^\text{[9][10][11]}\)

初等数论研究整数的一些方面，这些方面可以用初等方法来探讨，例如利用初等证明。\(^\text{[12]}\)相比之下，解析数论依赖于复数以及来自分析学和微积分的技巧。\(^\text{[13]}\)代数数论则使用代数结构（如域和环）来研究数的性质及其相互关系。几何数论运用几何学的概念研究数。\(^\text{[14]}\)数论的其他分支还包括：概率数论\(^\text{[15]}\)、组合数论\(^\text{[16]}\)、计算数论\(^\text{[17]}\)，以及应用数论，后者研究数论在科学与技术中的应用。\(^\text{[18]}\)
\subsubsection{起源}
\begin{figure}[ht]
\centering
\includegraphics[width=8cm]{./figures/b3c636666bde843e.png}
\caption{普林普顿 322 泥板文书} \label{fig_SHuLun_2}
\end{figure}
在有记录的历史中，数的知识存在于古代的美索不达米亚、埃及、中国和印度文明中\(^\text{[19]}\)。已知最早的算术性质的历史发现是普林普顿 322 号泥板，约公元前 1800 年。它是一块破碎的泥板，包含了一张毕达哥拉斯三元组（即整数$(a,b,c)$，满足$a^{2}+b^{2}=c^{2}$的数对）的列表。这些三元组数量众多且数值很大，不可能通过蛮力枚举得到\(^\text{[20]}\)。

该表的布局表明，它是借助于在现代语言中可写作如下的恒等式构造出来的\(^\text{[21]}\)：
$$
\left(\tfrac{1}{2}\left(x-\tfrac{1}{x}\right)\right)^{2}+1
=\left(\tfrac{1}{2}\left(x+\tfrac{1}{x}\right)\right)^{2},~
$$
这一恒等式隐含在古巴比伦日常的习题中\(^\text{[22]}\)。不过也有人提出，该表可能是作为学校习题的数值例子的来源\(^\text{[23][note 2]}\)。普林普顿 322 泥板是巴比伦数学中唯一现存的证据，可以被视为与现代所谓数论相对应的成果；然而，当时的一种巴比伦代数学要发达得多\(^\text{[24]}\)。

尽管其他文明可能在最初对希腊数学有所影响\(^\text{[25]}\)，但所有关于这种借鉴的证据都出现得相对较晚\(^\text{[26][27]}\)，因此很可能希腊的arithmētikḗ（数的理论性或哲学性研究）是一种本土传统\(^\text{[28]}\)。古希腊人对整除性产生了浓厚兴趣。毕达哥拉斯学派赋予完全数与亲和数神秘的性质。毕达哥拉斯的传统还谈及所谓的多边形数或图形数\(^\text{[29]}\)。欧几里得在其《几何原本》中专门论述了一些属于初等数论的主题，包括素数与整除性\(^\text{[30]}\)。他提出了用于计算两个数的最大公约数的欧几里得算法，并给出了一个证明，表明素数的数量是无限的。晚期古代最重要的arithmētikḗ权威是亚历山大的丢番图，他大约生活在公元三世纪。他撰写了《算术》，这是一部题集，其中的问题无一例外是寻找某个多项式方程组的有理数解，通常形式为：$f(x,y) = z^{2}$或$f(x,y,z) = w^{2}$.用现代术语来说，丢番图方程是指要求有理数或整数解的多项式方程。
\begin{figure}[ht]
\centering
\includegraphics[width=6cm]{./figures/759cf456e7ce5d29.png}
\caption{丢番图《算术》拉丁文译本（1621年，巴谢译）的书名页。} \label{fig_SHuLun_3}
\end{figure}
在罗马帝国灭亡之后，数论的发展转向亚洲，尽管过程断断续续。中国剩余定理出现在《孙子算经》（约三至五世纪）的一道习题中\(^\text{[31][32]}\)。这一结果后来被推广，并在秦九韶1247年的《数学九章》中以完整解法“大衍术”呈现\(^\text{[33][34]}\)。中国数学中也存在某些数秘学思想\(^\text{[note 3]}\)，但与毕达哥拉斯学派不同，这些思想似乎并未带来实质性的数学进展。尽管希腊天文学可能影响了印度的学术\(^\text{[35]}\)，但印度数学在其他方面似乎是一种本土传统\(^\text{[36][37]}\)。阿耶波多（Āryabhaṭa, 476–550 AD）提出，同时同余式对$n \equiv a_{1} \pmod{m_{1}}, \quad n \equiv a_{2} \pmod{m_{2}}$可以通过一种他称为kuṭṭaka（意为“粉碎法”或“逐次约化法”）的算法来求解【38】；这一过程与欧几里得算法十分接近\(^\text{[39]}\)。阿耶波多似乎主要考虑其在天文计算中的应用\(^\text{[35]}\)。婆罗摩笈多（Brahmagupta, 公元628年）则开启了对不定二次方程的系统研究，尤其是Pell方程。一种通解 Pell 方程的方法可能最早由Jayadeva发现；而最早保存下来的系统论述，则见于Bhāskara II（婆什迦罗二世）十二世纪的《种数论》\(^\text{[40]}\)。
\begin{figure}[ht]
\centering
\includegraphics[width=6cm]{./figures/3a3bcc555583aa3f.png}
\caption{西方视野中的伊本·海赛姆：在《月面学》的卷首插图中，‘Alhasen’（即伊本·海赛姆）代表理性之知，而伽利略则代表感官之知。} \label{fig_SHuLun_4}
\end{figure}
九世纪初，哈里发马蒙下令翻译大量希腊数学著作，以及至少一部梵文著作\(^\text{[41][42]}\)。丢番图的主要著作《算术》由古斯塔·伊本·卢卡（Qusta ibn Luqa, 820–912）译为阿拉伯文。卡拉吉（al-Karajī, 953 – 约1029）的部分著作《al-Fakhri》在一定程度上以此为基础。据Rashed Roshdi所言，卡拉吉的同时代人伊本·海赛姆已经知道后来被称为威尔逊定理的结果\(^\text{[43]}\)。在中世纪的西欧，除斐波那契所作关于等差数列平方数的论著外，几乎没有真正意义上的数论研究。直到文艺复兴晚期，随着对希腊古代著作的重新研究，这一局面才开始改变。其催化剂是丢番图《算术》的文本校订和译为拉丁文\(^\text{[44]}\)。
\subsubsection{近代早期数论}
\begin{figure}[ht]
\centering
\includegraphics[width=6cm]{./figures/c7a3992175aa4105.png}
\caption{皮埃尔·德·费马} \label{fig_SHuLun_5}
\end{figure}
法国数学家皮埃尔·德·费马（Pierre de Fermat, 1607–1665）从未正式发表过著作，而是通过通信以及在书页空白处的批注来交流思想\(^\text{[45]}\)。他对数论的贡献重新激发了欧洲对这一领域的兴趣。他提出了费马小定理（模算术中的一个基本结果）、费马大定理，并证明了费马直角三角形定理\(^\text{[2][46]}\)。他还研究了素数、四平方和定理以及Pell 方程\(^\text{[47][48]}\)。莱昂哈德·欧拉（Leonhard Euler, 1707–1783）对数论的兴趣最早可以追溯到 1729 年。当时，他的朋友、业余数学爱好者\(^\text{[note 4]}\)克里斯蒂安·哥德巴赫 向他推荐了费马在这一领域的一些工作\(^\text{[49][50]}\)。在费马未能成功引起同时代学者广泛关注之后，这一事件被称为现代数论的“再生”\(^\text{[51][52]}\)。欧拉证明了费马的一些断言，包括费马小定理；他率先开展了关于“每个整数都可以写成四个平方数之和”的证明工作\(^\text{[53]}\)；并处理了费马大定理的某些特例\(^\text{[54]}\)。他还研究了连分数与 Pell 方程之间的联系\(^\text{[55][56]}\)，并首次迈向了解析数论\(^\text{[57]}\)。
\begin{figure}[ht]
\centering
\includegraphics[width=6cm]{./figures/617824dcffca2fc8.png}
\caption{莱昂哈德·欧拉} \label{fig_SHuLun_6}
\end{figure}
三位欧洲同时代的数学家延续了初等数论的研究：约瑟夫–路易·拉格朗日（Joseph-Louis Lagrange, 1736–1813）给出了四平方和定理、威尔逊定理的完整证明，并发展了 Pell 方程的基本理论。阿德里安–玛丽·勒让德（Adrien-Marie Legendre, 1752–1833）提出了二次互反律，并猜想了相当于素数定理和狄利克雷算术级数定理的结果。他还系统研究了方程
$ ax^{2} + by^{2} + cz^{2} = 0$在他晚年时，他成为第一个证明费马大定理在 (n=5) 情形成立的人\(^\text{[59]}\)。卡尔·弗里德里希·高斯（Carl Friedrich Gauss, 1777–1855）在 1801 年撰写了《算术研究》，该书对数论产生了巨大影响，并为整个 19 世纪的研究奠定了议程。他在书中证明了二次互反律\(^\text{[60]}\)，并发展了二次型理论。他还为同余式引入了一些基本记号，并专门讨论了计算问题，包括素性检验\(^\text{[61]}\)。此外，他建立了单位根与数论之间的联系\(^\text{[62]}\)。通过这些工作，可以说高斯已在某种意义上为后来的埃瓦里斯特·伽罗瓦的研究，以及代数数论领域，开辟了道路。
\subsubsection{数论的成熟与分支划分}
\begin{figure}[ht]
\centering
\includegraphics[width=6cm]{./figures/f56c4c4e4969cc6c.png}
\caption{彼得·古斯塔夫·勒热纳·狄利克雷} \label{fig_SHuLun_7}
\end{figure}
自19世纪早期开始，数论逐渐经历了以下发展：
\begin{itemize}
\item 数论（或高等算术）作为一个独立研究领域的自觉兴起【63】。
\item 现代数学诸多工具的发展，这些工具是基础现代数论所必需的：复分析、群论、伽罗瓦理论——同时伴随着分析的更大 rigor（严格性）以及代数中的抽象化。
\item 数论粗略划分为现代分支，特别是解析数论与代数数论。
\end{itemize}
代数数论可说起始于对互反律与圆分问题的研究，但真正走向成熟则依赖于抽象代数、理想论与赋值论的早期发展（见后文）。解析数论的常规起点是狄利克雷算术级数定理（1837）【64】【65】，其证明引入了L-函数，并涉及一些渐近分析以及对实变量的极限过程【66】。

实际上，最早在数论中使用解析思想可追溯至欧拉（1730年代）【67】【68】，他使用了形式幂级数以及非严格（或隐含）的极限论证。将复分析引入数论则要更晚一些：黎曼（Bernhard Riemann, 1859）关于$\zeta$函数的研究是公认的起点【69】。而 雅可比的四平方和定理（1839）虽更早，却属于最初不同的研究方向，如今已在解析数论中占据核心地位（模形式）【70】。

美国数学会设立了柯尔数论奖。此外，数论还是费马奖所奖励的三大数学分支之一。
\subsection{主要分支}
\subsubsection{初等数论}
\begin{figure}[ht]
\centering
\includegraphics[width=6cm]{./figures/19f35a6cc5cac8ac.png}
\caption{数论学家保罗·埃尔德什与陶哲轩于1985年的合影，当时埃尔德什72岁，陶哲轩10岁} \label{fig_SHuLun_8}
\end{figure}
初等数论通过算术中的基本方法研究数论中的主题。[4] 它的主要研究对象包括整除性、因式分解和素数性，以及模算术中的同余关系。[71][12] 其他的研究内容还包括丢番图方程、连分数、整数拆分以及丢番图逼近。[72]

算术是研究数值运算的学科，探讨数字如何通过加法、减法、乘法、除法、乘方、开方和对数等算术运算进行组合和变换。例如，乘法是一种运算，它将两个数（称为因子）结合成一个数（称为积），如：$2 \times 3 = 6$【73】。

整除性是两个非零整数之间的一种性质，与除法相关。若一个整数$a$可以被一个非零整数 $b$整除，则称$a$是$b$的倍数；换言之，存在某个整数$q$，使得$a = bq$。另一种等价的表述是：$b$整除$a$，记作$b \mid a$。反之，若不满足整除关系，则$a$不能被$b$整除，除法会产生余数。欧几里得的除法引理断言，$a$和$b$一般可以表示为：$a = bq + r$,其中余数$r$是最小的正数剩余量。初等数论研究各种整除规则，以便快速判断一个给定整数是否可以被某个固定的除数整除。例如，众所周知，一个整数如果它的十进制各位数字之和能被 3 整除，那么该整数本身也能被 3 整除。[74][9][75]
\begin{figure}[ht]
\centering
\includegraphics[width=6cm]{./figures/13cc8818feba3d5f.png}
\caption{} \label{fig_SHuLun_9}
\end{figure}
若若干个非零整数有一个公约数，则该整数能够整除它们所有数。最大的这种约数称为最大公约数（gcd）。如果两个整数的最大公约数是 1，即它们唯一的公约数是 1，则称这两个整数互素或相对素数。欧几里得算法可用于计算两个整数 (a, b) 的最大公约数，其方法是反复应用除法引理，并在每一步将除数和余数交换位置。该算法还可以扩展来求解线性丢番图方程的特殊情况：$ax + by = 1$。一个丢番图方程一般具有多个未知数，并且系数都是整数。另一类丢番图方程体现在**勾股定理**中：$x^2 + y^2 = z^2$，如果其解都是整数，则称为勾股三元组。[9][10]另一种重要的表达形式是连分数，它把一个数写成一个整数与一个分数之和，而分母中的分数又是类似形式的整数与分数之和。[76]

初等数论研究整数的整除性质，例如奇偶性（偶数与奇数）、素数和完全数。一些重要的数论函数包括：约数个数函数、约数求和函数及其变形，以及欧拉函数。素数是指大于 1 的整数，且它的正约数只有 1 和它自身。大于 1 的正整数若不是素数，则称为合数。欧几里得定理证明了素数有无穷多个，即集合${2, 3, 5, 7, 11, \dots}$。埃拉托色尼筛法是一种高效算法，它通过逐步剔除所有合数来识别不超过给定自然数的所有素数。[77]

因数分解是一种将一个数表示为乘积的方法。特别是在数论中，整数分解是将一个整数分解为若干整数的乘积。不断重复这一过程直到所有因子都是素数，被称为素因数分解。素数的一个基本性质体现在**欧几里得引理**中。该引理的推论是：如果一个素数能整除若干整数的乘积，那么它必然能整除其中至少一个因子。与素因数分解相关的最基本定理是算术基本定理（又称唯一分解定理）。该定理指出：每个大于 1 的整数都可以分解为素数的乘积，并且这种分解在因子顺序之外是唯一的。例如：$120 = 2 \times 2 \times 2 \times 3 \times 5 = 2^{3} \times 3 \times 5$。[78][9]

模算术处理有限集合中的整数，并引入同余和剩余类的概念。若两个整数$a, b$对模$n$（一个正整数，称为模数）满足$n \mid (a-b)$，则称$a$与$b$在模$n$下同余。即在对$a$与$n$、$b$与$n$进行欧几里得除法时得到相同的余数。这写作：$a \equiv b \pmod{n}$。类似于 12 小时制时钟，$4 + 9 = 13$，但在模 12 的意义下同余于 1。一个模$n$的剩余类是包含所有与某个指定整数$r$在模$n$下同余的整数的集合。例如，$6\mathbb{Z} + 1$包含所有能表示为$6k+1$的整数。

模算术提供了许多公式，可以快速解决涉及大指数幂的同余问题。其中一个重要定理是费马小定理，它指出：如果素数$p$与某个整数$a$互素，则$a^{p-1} \equiv 1 \pmod{p}$。欧拉定理将其推广为：对于任意整数$n$，若$a$与$n$互素，则
$$
a^{\varphi(n)} \equiv 1 \pmod{n}~
$$
其中$\varphi(n)$是欧拉函数，表示小于等于$n$且与$n$互素的正整数的个数。模算术还提供了解决含有未知数的同余方程的公式，类似于代数方程的解法，例如中国剩余定理。[79]
\subsubsection{解析数论}
\begin{figure}[ht]
\centering
\includegraphics[width=6cm]{./figures/929b80c6eb74e539.png}
\caption{复平面上的黎曼$\zeta$函数$\zeta(s)$。点$s$的颜色表示 =$zeta(s)$ =的取值：深色表示值接近零，而色相表示该值的辐角。} \label{fig_SHuLun_10}
\end{figure}
与初等数论相对，解析数论依赖于复数以及分析与微积分的技巧。解析数论可以从两个角度来定义：
\begin{itemize}
\item 从方法角度：它是利用实分析和复分析工具研究整数的学科；[64]
\item 从研究对象角度：它是数论中关于某些数（例如素数）的大小和密度的估计的研究，而不是单纯关于恒等式的研究。[80]
\end{itemize}
解析数论研究的内容包括：素数的分布、数论函数的行为以及无理数等问题。[81]
\begin{figure}[ht]
\centering
\includegraphics[width=8cm]{./figures/35f5b602fc3a06cc.png}
\caption{模群在上半平面的作用。灰色区域是标准基本域。} \label{fig_SHuLun_11}
\end{figure}
数论素有“结论浅显，证明深奥”的名声。许多数论结论可以用浅显的方式告诉普通人，但它们的证明往往并不容易理解，部分原因在于所使用的工具范围极其广泛，在数学中也算是罕见的。[82]解析数论中的经典问题包括：素数定理、哥德巴赫猜想、孪生素数猜想、哈代–李特伍德猜想、华林问题以及黎曼猜想。解析数论的重要工具有：圆方法，筛法，L-函数（以及对其性质的研究）。此外，模形式理论（以及更一般的自守形式理论）在解析数论工具箱中的地位也日益核心。[83]

分析学是数学的一个分支，研究极限——即当自变量（或序列的指标）趋近某个特定值时，序列或函数所趋近的数值。例如，数列 0.9, 0.99, 0.999, … 的极限是 1。在函数的语境中，当$x \to \infty$时，函数$\frac{1}{x}$的极限是 0。[84]复数是实数的扩展，引入了虚数单位$i$，其定义满足$i^2 = -1$。【85】 每一个复数都可以表示为$x + i y$，其中
$x$称为实部$y$称为虚部。

素数的分布由函数$\pi(x)$描述，该函数统计所有小于等于给定实数的素数个数。素数的分布是不可预测的，是数论中的主要研究对象之一。虽然已经发展出一些素数的初等公式，例如欧拉的素数多项式，能生成部分素数序列，但随着素数变大，这些公式就不再适用。解析数论中的素数定理形式化了这样一种观念：素数在数值增大的同时出现得越来越稀疏。一个经典的近似分布陈述是：$\frac{x}{\log(x)}$能够很好地逼近$\pi(x)$。另一种更精确的分布形式涉及到偏移的对数积分函数，它比$\frac{x}{\log(x)}$更快地收敛到$\pi(x)$。[3]
\begin{figure}[ht]
\centering
\includegraphics[width=8cm]{./figures/392525ce78a681e7.png}
\caption{利用$\zeta$函数的零点对素数计数函数估计的修正} \label{fig_SHuLun_12}
\end{figure}
$\zeta$函数已被证明与素数的分布密切相关。它被定义为级数：
$$
\zeta(s) = \sum_{n=1}^{\infty}\frac{1}{n^s} = \frac{1}{1^s} + \frac{1}{2^s} + \frac{1}{3^s} + \cdots~
$$
当$s > 1$时，该级数收敛。欧拉证明了它与所有素数的无穷乘积之间存在联系，给出了如下恒等式：
$$
\zeta(s) = \prod_{p \ \text{prime}}\left(1 - \frac{1}{p^s}\right)^{-1}.~
$$
黎曼将$\zeta$函数的定义推广到了复变量，并提出猜想：在所有非平凡零点（即$0 < \Re(s) < 1$区间内使$\zeta(s)=0$ 的解）中，$s$的实部都等于$\frac{1}{2}$.他建立了这些非平凡零点与素数计数函数之间的联系。这就是如今著名的黎曼猜想，它仍然是数学中的悬而未决的问题。如果该猜想得到解决，将会对理解素数的分布带来直接而深远的影响。[86]

人们可以针对代数数提出分析性的问题，并利用分析方法来解答这些问题；正是在这一点上，代数数论与解析数论相交。例如，可以定义素理想（在代数数域中对素数的推广），并询问在某个范围内共有多少个素理想。这个问题可以通过研究德德金$\zeta$函数来回答，该函数是黎曼 $\zeta$函数的推广，而黎曼$\zeta$函数则是该领域最核心的解析对象之一。[87]这就是解析数论中的一种普遍方法：通过研究某个适当构造的复值函数的解析性质，来推导关于某个数列（此处是素理想或素数）分布的信息。[88]

初等数论依靠的是所谓的“初等证明”，该术语排除了复数的使用，但可能会包括一些基础分析方法。[72]例如，素数定理在 1896 年首先是用复分析证明的，但直到1949年，Erdős和Selberg才找到了一个“初等证明”。[89]不过，“初等”这一术语带有一定的模糊性。例如，基于复数Tauberian 定理（如 Wiener–Ikehara 定理）的证明，通常被认为非常有启发性，但尽管它们只使用了傅里叶分析而非复分析，却通常不被视为“初等证明”。在这里和其他情境中，一个“初等证明”可能比更高深的方法对大多数读者来说更冗长、更难以理解。

一些通常被视为解析数论组成部分的主题（例如筛法理论），若以“工具”为标准，可能更适合归入第二种定义而非第一种。[note 5]例如，小筛法几乎不涉及分析手段，但仍然属于解析数论的范畴。[note 6]
\subsubsection{代数数论}
代数数是指满足某个有理系数多项式方程$f(x)=0$的复数。例如，方程$x^{5}+\tfrac{11}{2}x^{3}-7x^{2}+9=0$的每个解$x$都是代数数。由代数数组成的域称为代数数域，或简称数域。代数数论主要研究代数数域。[90]

可以认为，最简单的一类数域，即二次扩域，已经被高斯研究过了，因为他在《算术研究》中关于二次型的讨论，可以用现代语言改写为关于二次域中的理想与范数的讨论。（一个二次域由所有形如$a+b\sqrt{d}$的数构成，其中$a, b$是有理数，而$d$是一个平方根不是有理数的有理数。）事实上，十一世纪的chakravala方法，用现代术语来说，本质上就是一个求解实二次数域的单位元的算法。然而，无论是Bhāskara还是Gauss，在他们的时代都尚未意识到“数域”这一概念本身。

这一学科的基础在19世纪晚期得以奠定，当时引入了理想数、理想理论和赋值理论，这三者是互补的方法，用于处理代数数域中唯一分解失效的问题。（例如，在由有理数与$\sqrt{-5}$生成的数域中，数$6$可以分解为$6 = 2 \cdot 3$也可以分解为$6 = (1+\sqrt{-5})(1-\sqrt{-5})$,其中$2, 3, 1+\sqrt{-5}, 1-\sqrt{-5}$都是不可约的，因此在朴素意义上类似于整数中的素数。）理想数（由库默尔 Kummer 引入）的发展最初似乎受到高次互反律研究的推动，[91] 即二次互反律的推广。

数域通常作为较小数域的扩张来研究：如果域$L$含有域$K$，则称$L$是$K$的扩张。（例如，复数域$\mathbf{C}$是实数域$\mathbf{R}$的扩张，而实数域$\mathbf{R}$ 又是有理数域$\mathbf{Q}$的扩张。）对一个给定数域的所有可能扩张进行分类，是一个困难且部分未解决的问题。阿贝尔扩张（即$L/K$的伽罗瓦群$\mathrm{Gal}(L/K)$是阿贝尔群的扩张）相对比较容易理解。它们的分类是类域论的核心目标，这一研究计划由Kronecker和 Eisenstein在 19 世纪晚期部分发起，并在1900–1950年间得到大规模发展和完成。

代数数论的一个活跃研究领域是岩泽理论。而当前数学中最重要的大规模研究计划之一——朗兰兹纲领，有时被描述为一种尝试：将类域论推广到非阿贝尔数域扩张。
\subsubsection{丢番图几何}
丢番图几何的核心问题是：确定一个丢番图方程在何种情况下有整数或有理数解；如果存在，那么解有多少。其研究方法是将方程的解视作一个几何对象。

例如，两个变量的方程定义了平面上的一条曲线。更一般地，两个或多个变量的方程或方程组，可以定义$n$维空间中的一条曲线、一个曲面或其他几何对象。在丢番图几何中，人们关心的问题是：在这些曲线或曲面上是否存在有理点（所有坐标都是有理数的点）或整点（所有坐标都是整数的点）。如果存在这样的点，下一步则是问：它们有多少个？分布情况如何？一个基本的问题是：在给定的曲线或曲面上，是否存在有限个还是无限多个有理点。

例如，考虑毕达哥拉斯方程：$x^{2} + y^{2} = 1$我们想知道它的有理解，即满足$x$和$y$都是有理数的点$(x,y)$。这与寻找所有整数解$a^{2} + b^{2} = c^{2}$是等价的；任何一个整数解都会给出一个有理解：$x = a/c, \quad y = b/c$.换句话说，这就是在曲线$x^{2} + y^{2} = 1$上寻找所有有理坐标点。该曲线是以原点为圆心、半径为 1 的圆。
\begin{figure}[ht]
\centering
\includegraphics[width=8cm]{./figures/1a10c45a3990694f.png}
\caption{两个椭圆曲线的例子，即亏格为 1 且至少有一个有理点的曲线。} \label{fig_SHuLun_13}
\end{figure}
将关于方程的问题重新表述为关于曲线上点的问题，是一种非常巧妙的方式。代数曲线（即两个变量的多项式方程$f(x,y)=0$）上有理点或整点的数量是否有限，关键取决于该曲线的亏格。[note 8] 这一研究方法的一项重大成果就是怀尔斯对费马大定理的证明，其中其他几何学概念同样发挥了至关重要的作用。

另一个紧密相关的领域是丢番图逼近：给定一个数$x$，研究它能够被有理数逼近的精确程度。人们寻找相对于分母书写复杂度而言“良好”的逼近。具体来说，如果$\gcd(a,q)=1, \quad \left|x-a/q\right| < \frac{1}{q^{c}}$,其中$c$较大，则称 $\tfrac{a}{q}$是$x$的一个良好逼近。这个问题在$x$为代数数时尤其重要。如果某个数不能被很好地逼近，那么某些方程就不会有整数或有理解。此外，若干概念（尤其是“高度”height）在丢番图几何和丢番图逼近中都至关重要。这一问题在超越数论中也具有特殊意义：如果一个数能被比任何代数数更好地逼近，那么它就是一个超越数。正是通过这种论证，人们证明了$\pi$和$e$是超越数。

需要注意的是，丢番图几何不应与数的几何混淆，后者是一种利用图解方法回答某些代数数论问题的工具集合。而算术几何则是丢番图几何的当代理论名称，尤其当人们希望强调它与现代代数几何（如法尔廷斯定理）之间的联系时，而不是仅仅强调与丢番图逼近技术的关联。
\subsubsection{其他分支领域}
概率数论从这样的问题开始：随机取一个 1 到一百万之间的整数$n$。它是质数的可能性有多大？（这其实就是在问一百万以内有多少个质数）。平均而言，$n$会有多少个质因子？它比平均情况多得多或少得多的因子或质因子的概率是多少？

数论中的组合学从这样的问题开始：一个相当‘稠密’的无限集合$A$是否包含许多处于等差数列中的元素：

$a, , a+b, , a+2b, , a+3b, \ldots, a+10b$？是否有可能把大的整数写成集合$A$中元素的和？”
\begin{figure}[ht]
\centering
\includegraphics[width=6cm]{./figures/9b65f336e2f2d3c7.png}
\caption{莱默筛法机，一种原始的数字计算机，用于寻找质数并解简单的丢番图方程。} \label{fig_SHuLun_14}
\end{figure}
这里有两个主要问题：‘能否计算？’以及‘能否快速计算？’任何人都可以检验一个数是否为质数，或者在它不是质数时将其分解为质因数；但要快速完成则是另一回事。如今已经有了快速的质数判定算法，但尽管在理论和实践上付出了大量努力，目前仍没有真正快速的因数分解算法。
\subsection{应用}
长期以来，数论整体，尤其是素数研究，一直被视为纯数学的典范例子，在数学之外几乎没有任何应用，除了利用素数齿轮来均匀分配磨损这一点之外。[92] 尤其是像英国数学家 G. H. Hardy 这样的数论学家，以从事完全没有军事意义的工作为荣。[93] 数论学家 Leonard Dickson（1874–1954）甚至说过：“感谢上帝，数论未被任何应用玷污。”然而，这种观点如今已不再适用于数论。[94]

这种关于数论纯粹性的想法在 1970 年代被彻底打破，因为当时公开宣布素数可以用作公钥密码算法的基础。[95] 像 RSA 这样的方案就建立在将大型合数分解为素因子的困难性之上。[96] 这些应用促使人们对处理素数的算法进行了深入研究，尤其是素性测试方法，即判断给定数是否为素数的手段。素数还被用于计算机科学中的校验和、哈希表以及伪随机数生成器。

1974 年，Donald Knuth 说过：“几乎所有初等数论中的定理，都以一种自然、合理的方式出现，并且与让计算机进行高速数值计算的问题紧密相关。”[97] 初等数论常在计算机科学的离散数学课程中教授，它在数值分析中的连续问题上也有应用。[98]

如今，数论已经在多个现代领域中找到了应用，涉及范围广泛，例如：
\begin{itemize}
\item 计算机科学：快速傅里叶变换（FFT）算法可高效计算离散傅里叶变换，在信号处理和数据分析中有重要应用。[99]
\item 物理学：黎曼猜想与素数的分布相关，并因其在物理学中的潜在意义而受到研究。[100]
\item 纠错码：有限域理论和代数几何被用于构造高效的纠错码。[101]
\item 通信：蜂窝移动网络的设计需要用到模形式理论，而模形式是解析数论的一部分。[102]
\item 音乐音阶研究：所谓“十二平均律”概念，是大多数现代西方音乐的基础，其思想是将一个八度平均分成 12 个等份。[103] 这一思想已经通过数论，尤其是 2 的 12 次方根的性质得到研究。
\end{itemize}
\subsection{另见}
\begin{itemize}
\item 算术动力学
\item 代数函数域
\item 算术拓扑
\item 有限域
\item p 进数
\item 数论算法列表
\end{itemize}
\subsection{注释}
\begin{enumerate}
\item “Arithmetic”（算术）一词可能重新获得了一定的地位，可以说是受法国学界的影响。例如参见 Serre 1996。一九五二年，Davenport 仍需特别说明他所指的是 The Higher Arithmetic。Hardy 和 Wright 在其《An Introduction to the Theory of Numbers》（1938）的引言中写道：\\
“我们曾一度打算将书名改为 An Introduction to Arithmetic，这是一个更为新颖、在某些方面也更为恰当的标题；但有人指出，这可能会导致对书的内容产生误解。”（见 Hardy & Wright 2008）
\item Robson 2001，第 201 页。这一问题具有争议性，参见 Plimpton 322。Robson 的文章以一种带有论战性的笔法写成（Robson 2001，第 202 页），其目的在于“也许……把 [Plimpton 322] 从神坛上拉下来”（Robson 2001，第 167 页）；但同时，他的结论是：\\
“……问题‘这块泥板是如何计算出来的？’的答案，不必与‘这块泥板提出了什么问题？’的答案相同。第一个问题，最令人满意的答案是倒数对，这在半个世纪前就有人提出过；而第二个问题，则可能是某种直角三角形问题。”（Robson 2001，第 202 页）\\
Robson 质疑了这样一种看法：即认为制作Plimpton 322的书吏（他必须“为生计而工作”，不可能属于“悠闲的中产阶级”）会在缺乏“新数学市场”的情况下，单凭自己“无聊的好奇心”去完成这一工作。（Robson 2001，第 199–200 页）
\item 例如，参见《孙子算经》第三章第 36 题（引自 Lam & Ang 2004，第 223–224 页）：\\
[36] 题现有一孕妇，年二十九岁。若妊娠期为九个月，求胎儿性别。答曰：男。
术：下四十九，加妊娠月数，减去其年。余数再依次减去：一为天，二为地，三为人，四为四时，五为五行，六为六律，七为七星（北斗），八为八风，九为九州。若余数为奇，则为男；若余数为偶，则为女。\\
这是《孙子算经》中本来简明实用的论述中出现的最后一个问题。
\item 直到十七世纪下半叶之前，学术职位仍然非常稀少，大多数数学家和科学家必须通过其他方式谋生（Weil 1984，第 159、161 页）。 （当时已经出现了一些可辨识的职业实践特征，例如：寻找通信伙伴、拜访国外同事、建立私人图书馆（Weil 1984，第 160–161 页）。）局势在十七世纪晚期开始发生转变（Weil 1984，第 161 页）；科学学会在英格兰（皇家学会，1662 年）、法国（科学院，1666 年）和俄国（1724 年）相继成立。Euler 于 1726 年受邀前往俄国科学院，并于 1727 年抵达圣彼得堡（Weil 1984，第 163 页；Varadarajan 2006，第 7 页）。在这一背景下，“业余爱好者”这一通常用于形容 Goldbach 的术语有其明确含义并且颇为合理：他被描述为一位靠做间谍谋生的文人（Truesdell 1984，第 xv 页；引自 Varadarajan 2006，第 9 页）。不过需要注意的是，Goldbach 也发表过一些数学著作，并且有时还曾担任过学术职位。
\item 筛法理论在许多标准论著中被视为解析数论的主要分支之一；例如，参见 Iwaniec & Kowalski 2004 或 Montgomery & Vaughan 2007。
\item 这一说法适用于某些组合型筛法，例如 Brun 筛，而不适用于“大筛”。对后者的研究现已包括调和分析和泛函分析的思想。
\item 扩域$L/K$的伽罗瓦群由那些把$L$中元素映射到$L$中其他元素、同时保持 (K) 中所有元素不变的运算（同构）组成。比如，$\mathrm{Gal}(\mathbb{C}/\mathbb{R})$含有两个元素：恒等映射（把复数$x+iy$映射到自身）以及复共轭（把 $x+iy$映射到$x-iy$的映射）。扩域的伽罗瓦群能够揭示许多关于该扩域的关键性质。
对伽罗瓦群的研究始于Évariste Galois。用现代语言表述，他工作的主要成果是：一个方程$f(x) = 0$可以用根式解（即可以用四则运算以及平方根、立方根等表达$x$）当且仅当由该方程的根对有理数域所作的扩域，其伽罗瓦群在群论意义下是可解的。这里的“可解”是群论中的一个简单性质，对于有限群可以很容易检验。
\item 亏格可以这样定义：允许方程$f(x,y) = 0$中的变量取复数，那么该方程定义了一个嵌入在（射影的）四维空间中的二维曲面（因为两个复数变量可以分解成四个实数变量，即四个维度）。该曲面上甜甜圈状“洞”的数量称为曲线$f(x,y) = 0$的亏格。
\end{enumerate}
\subsection{参考文献}
\begin{enumerate}
\item Long 1972, 第1页。
\item Karatsuba, A.A. (2020). 《数论》 Encyclopedia of Mathematics. Springer. 检索日期：2025-05-03。
\item Moore, Patrick (2004). “Number theory”。见：Lerner, K. Lee；Lerner, Brenda Wilmoth（编）. The Gale Encyclopedia of Science. 第4卷（第3版）. Gale. ISBN 0-7876-7559-8。
\item Tanton, James (2005). “Number theory”。Encyclopedia of Mathematics. 纽约：Facts On File. 第359–360页. ISBN 0-8160-5124-0。
\item 洛萨诺-罗夫莱多 2019，第 xiii 页\\
\item 纳格尔与纽曼 2008，第 4 页\\
罗曼诺夫斯基 2008，第302–303页\\
HC 编辑部 2022b\\
MW 编辑部 2023\\
布赫什塔布与佩查耶夫 2020

\end{enumerate}